\section{Reward extraction across successor representations}
\label{app:reward-extraction}

\cref{tab:manipulation} gives every method the task as a goal observation, which is the interface the goal-conditioned reference implementations are built for.
This appendix gives the same frozen checkpoints the task as a \emph{reward} instead: each method draws $10^4$ transitions from its own training data, labels them with the task reward, and solves for the conditioning vector whose implied values best explain those labels, through its own inference rule.
Only successor representations appear here, because the comparison is about the reward-extraction step and the goal-conditioned baselines have none.

\begin{table}[h]
    \centering
    \small
    \caption{\textbf{Task given as a reward} (success rate, \%). The same $10^6$-step checkpoints as \cref{tab:manipulation}; the conditioning vector is inferred from $10^4$ reward-labelled dataset states, five tasks per environment, ten episodes per task, mean $\pm$ standard deviation over three seeds.}
    \label{tab:reward-extraction}
    \begin{tabular}{lrrrrr}
    \toprule
    Task & FB & HILP & RLDP & \name{} & \name{}$^{\,\mathrm{oc}}$ \\
    \midrule
    \texttt{cube-single} & 42.0\pmse{11.4} & 72.7\pmse{3.4} & 36.7\pmse{2.5} & \textbf{92.0}\pmse{0.0} & 84.7\pmse{4.1} \\
    \texttt{cube-double} &  4.7\pmse{0.9} &  8.0\pmse{0.0} &  2.0\pmse{1.6} & 14.7\pmse{4.1} & \textbf{34.0}\pmse{8.6} \\
    \texttt{scene}       & 10.7\pmse{6.8} & 37.3\pmse{4.1} & 10.7\pmse{4.1} & 50.7\pmse{6.8} & \textbf{52.0}\pmse{1.6} \\
    \texttt{puzzle-3x3}  &  4.0\pmse{3.3} &  0.7\pmse{0.9} & \textbf{9.3}\pmse{7.5} &  5.3\pmse{2.5} & 6.7\pmse{1.9} \\
    \bottomrule
    \end{tabular}
    \par\vspace{2pt}
    \begin{minipage}{\linewidth}\scriptsize\raggedright
    $^{\mathrm{oc}}$: the backward network reads the object-coordinate block of the observation instead of the full state. Every other column reads the full state.
    On \texttt{cube-double} reward extraction is capped near half of goal-conditioned success for every method: the environment re-permutes the cube-to-target assignment at every reset, while the reward samples encode one fixed assignment.
    \end{minipage}
\end{table}

Two readings follow.

First, the reward-extraction step is where the successor representations separate.
\name{} leads on three of the four environments, and its margin over FB is larger here than under goal conditioning ($92.0$ against $42.0$ on \texttt{cube-single}, $50.7$ against $10.7$ on \texttt{scene}).

Second, the two interfaces are not interchangeable, and which one a method prefers is itself informative.
FB and RLDP lose most of their performance when the task arrives as a goal rather than a reward ($42.0\to17.3$ and $36.7\to10.7$ on \texttt{cube-single}): evaluating their backward map at a single state is a harder query than averaging it against $10^4$ reward-weighted samples.
HILP moves the other way ($72.7\to86.0$), since its features are trained to make state-to-state distances meaningful and a goal is its native query.
\name{} is stable across both ($92.0$ and $84.3$), which is what a normalized successor measure predicts: the value is an average of the reward against a distribution, so it does not depend on whether the reward is concentrated at one state or spread over many.
